\begin{center}
{\Large\bfseries Descent of a residual chain}\\[4pt]
{\large The spectral frame inequality and the rank-three count}\\[4pt]
Edition 0.61 \quad 24 November 2026
\end{center}

A rank-four residual representation in the Dirichlet-GRH argument is
joined to a seed by a finite chain of polynomial deformations.
To use the chain at a finite level one must descend the deformations,
the local data and the common endpoints to the same residue field.
The first chapter proves that one coefficient ring works for everything
and for all subsequent enlargements. This proves the entire
\texttt{CHAIN\_DESC4} obligation.

The next two chapters bound the prime attached to every late word
and transfer an absolutely convergent operator expansion back to the
original spectral datum. The last part follows a rank-three counting
argument from its arithmetic host to the precise remaining
deformation-theoretic inputs.

\begingroup\emergencystretch=2em
\section{One finite coefficient field for a rank-four residual chain}
Fix a prime $\ell$, a coefficient ring $O_c$ in the prescribed algebraic
closure of $\mathbb Q_\ell$, and the valuation ring $V$ in the
algebraically closed field of the given ambient datum. Let $k$ be its residue field. For each admissible $O_d$ write
$k_d$ for its residue field and $j_d:k_d\hookrightarrow k$ for the map
induced by the given local map $O_d\to V$. These maps are compatible
under inclusions of coefficient rings. The ambient hypothesis says that
each element of $V$ agrees modulo its maximal ideal with the image of
some element of some $O_d$.

Let $\Gamma$ be the supplied compact space with group structure and
$\Phi$ its given frame over $O_c$. Write $I$ for the finite set of local
labels. Over a ring $A$, a point $P=(r,n)$ consists of a representation
$r:\Gamma\to M_4(A)$ and one matrix $n_v$ for each $v\in I$. A link
from $P_0$ to $P_1$ is a point over $A[X]$ satisfying the original
residual conditions, with evaluations $P_0$ at zero and $P_1$ at one.
A rank-four chain is a good initial point followed by a finite sequence
of links with good terminal points. Goodness here means the disjunction
$\mathrm{Good4}$ specified below. The polynomial links themselves must
descend.

% BEGIN RCD1
\paragraph{RCD1. A uniform ring for all entries of a finite list of points.}
For a finite list of points over $k$ or $k[X]$ whose representation
entries are locally constant, there is one admissible $O_d$ such that,
for every $O_{c'}\supseteq O_d$, every point in the list has a unique
preimage over $k_{c'}$ or $k_{c'}[X]$, respectively. We may take
$O_d\supseteq O_c$.
\begin{proof}
For $z\in k$, choose a lift to $V$ and use the ambient hypothesis to
find $y\in O_b$ with residue image $z$. For every $O_{c'}\supseteq O_b$,
the residue of the same $y$ in $k_{c'}$ maps to $z$.
Any two coefficient rings lie in a third. Their fraction fields are
finite extensions of $\mathbb Q_\ell$ inside its fixed algebraic
closure. The ring of integers of their finite compositum is an
admissible coefficient ring and contains both by integrality.
Induction gives a common ring for any finite list of residues and
the initial $O_c$.

For each entry of $r$, the open level sets of its locally constant
value cover $\Gamma$. A finite subcover gives finitely many values.
There are sixteen entries and only finitely many label matrices.
In the polynomial case this first gives finitely many \emph{polynomials};
the union of their coefficient sets is finite. Apply the preceding
argument once to all the coefficients, across all the points.

For any larger $O_{c'}$, both $j_{c'}$ and its coefficientwise polynomial
extension are injective. Take the unique preimage of each entry.
Matrix multiplication commutes with these inclusions. Injectivity
reflects $r(1)=1$ and $r(gh)=r(g)r(h)$, giving a representation over
the smaller ring. Together with the preimages of the $n_v$ it gives the
unique point. Its representation entries remain locally constant
because their level sets are the same.
\end{proof}
% END RCD1

% BEGIN RCD2
\paragraph{RCD2. Residual conditions and good vertices descend.}
Let $f:A\hookrightarrow B$ be an injective unital map of commutative
rings, let $a:O_c\to A$ be the coefficient map and let $P$ be an
$A$-point on $\Phi$. If its image satisfies $\Phi.\mathrm{IsRes}$
over $f a$, then $P$ satisfies it over $a$. The same implication holds
for $\mathrm{Good4}$.
\begin{proof}
The symplectic adjoint is given by transposition and fixed signed
changes of indices, so it commutes with $f$. Matrix multiplication,
the similitude equation and every entry of
$\operatorname{adj}(n_v)r(g)n_v$ commute with $f$. In particular so
do the diagonal functions $\psi_{v,i}(g)$.

The pair must have the prescribed multiplier, kill $\Phi.N$, be diagonal
at the distinguished element, and have lower unitriangular symplectic
$n_v$. On $\Phi.G_v$ the conjugated matrices must be upper triangular,
and on $\Phi.\mathrm{Mu}_v$ all their diagonal entries must be one.
The residual data prescribe the diagonals of $r(g)$ and of the local
conjugates by the characters $a(\chi_i)$. They also require one of
two support patterns: global support in $n_P$ together with $n_v$
supported in $q_Q$ at every label, or global support in $n_M$.
Each support set specifies the permitted off-diagonal weights.
Every equality and required zero follows by injectivity, keeping the
given support branch. The continuity clause means continuity into the
discrete quotients by powers of the zero ideal, entry by entry.
Exponent one gives local constancy, which an injection reflects.
Conversely every quotient of a locally constant entry is locally
constant. This proves the full residual predicate.

For $\mathrm{Good4}$, a specified root occurs when some $g,i,j$ of that
weight have $r(g)_{ij}\ne0$. Keep those witnesses. The two extra branches
require global support in $a_A$ and $(n_v)_{10}\ne0$ at each label,
together with their specified occurring roots. The ordinary good
branches use occurrences of roots alone or the condition $\mathrm{InV}$,
which requires global support in $v_V$, local support in $q_Q$, and
$(n_v)_{20}\ne0$. Its last branch also chooses, at each label, an
$s\in\Phi.G_v$ where a specified root character differs from one.
Keep the branch and these same $s$. A zero or equality to one
downstairs would give that equality upstairs. This exhausts the tests.
\end{proof}
% END RCD2

% BEGIN RCD3
\paragraph{RCD3. Descent of the entire directed chain.}
Suppose $s,t$ are points over $k$ with
$\Phi.\mathrm{IsRes}(\mathrm{red}_c,s)$ and
$\mathrm{Chain4}(\Phi,\mathrm{red}_c,s,t)$. There is a coefficient
datum $d$ such that for every $c'$ with $O_c,O_d\subseteq O_{c'}$
there are $s',t'$ over $k_{c'}$ mapping to $s,t$. Over the
coefficient map $a:O_c\hookrightarrow O_{c'}\twoheadrightarrow k_{c'}$,
$s'$ satisfies $\Phi.\mathrm{IsRes}(a,s')$ and
$\mathrm{Chain4}(\Phi,a,s',t')$ holds. This is the original
$\mathrm{ChainDesc4}$.
\begin{proof}
Choose the finite link witnesses $P_1,\ldots,P_m$ for the given
reflexive transitive closure. The residual hypothesis on $s$ and
the residual hypothesis on each \emph{polynomial} $P_j$ give local
constancy of their representation entries by the zero-ideal test in
RCD2. Apply RCD1 to this finite list including $s$. For each allowed
$c'$ let $s'$ and $Q_j$ be their unique preimages. On constants
$j_{c'}a=\mathrm{red}_c$ by compatibility of the maps into $V$.
The corresponding square also commutes after including constants
into the polynomial rings. RCD2 gives the residual predicate for
$s'$ and for every $Q_j$ over its polynomial coefficient map.

Evaluation at zero and one commute with coefficientwise inclusion.
Thus $Q_1(0)=s'$ and $Q_j(1)=Q_{j+1}(0)$ whenever the indicated
links exist. In each case the two points agree after an injection
on both $r$ and every $n_v$. Take $t'=Q_m(1)$ for $m>0$.
RCD2 gives goodness of $s'$ and of each $Q_j(1)$ from the original
witnesses. This is a chain with exactly the original orientation
and length. For $m=0$ take $t'=s'$ and use initial goodness.
The ring $O_d$ was fixed before $c'$, giving the claimed quantifiers.

The Case-III application can therefore enlarge coefficients again
to include a seed and other integral data while retaining the chain
at the final ring. Propagating a condition on components of the
deformation space along it uses the separate rank-four theorem.
\end{proof}
% END RCD3

\section{The frame inequality at the selected level prime}

Write $T_u=(D.\mathrm{mem}(u)).A.T$ for the prime already assigned
to the word $u$. Let $g_0,g_1$ be the two host primes, $v$ the
host root discriminant divided by two, $m$ the number of bottom
weights. The rank record requires even $m\ge2$ and defines
$r_S=m/2-1$, so $2r_S=m-2$ also for its natural-number arithmetic.
The scale in the frame comparison is
$\mathrm{kaS}(D)=\pi/(v g_0 g_1)$. Thus a lower bound on $T_u$
has to retain the dependence on \emph{both} host primes and on
$v$. The arithmetic window already contains these three factors.

% BEGIN TFR1
\paragraph{TFR1. The lower edge of the prime window with its host factors.}
\label{claim:TFR1}
Use the ranks $m,k,f,r_S$ of an original datum and set
\[
 A=\log(vg_0g_1/\pi),\qquad
 X_u=\frac{1}{fk}\log\frac{D_*}{d_*(u)},\qquad L=x^\epsilon.
\]
Use the hypotheses of QW6 on the original row and column at $u$,
including $k\le x^{C_*}$, and the positive prime-window data
$T_u\ge (D_*/d_*(u))^{1/(fk)}\Gamma_u^{-2/(fk)}$.
If $m^2\le k$, $L\ge1$ and all the $m$ middle weights at every
place satisfy $w_{j\rho}\ge e$, then
\[
 \log T_u\ \ge\ mA+m-2C_\epsilon L
                    -\frac{2\log\Gamma_u}{fk},
 \qquad C_\epsilon=30+\frac{2+2C_*}{\epsilon}.       \tag{TFR.1}
\]
In particular, if $m\ge 2C_\epsilon L+2$ and
$2\log\Gamma_u/(fk)\le2$, then
\[
 T_u\ge(vg_0g_1/\pi)^m,\qquad
 T_u\,\mathrm{kaS}(D)^{2r_S}\ge(vg_0g_1/\pi)^2\ge1.
                                                        \tag{TFR.2}
\]

\begin{proof}
The weights $w_{j\rho}$ are the interlacing entries called $e_j$
in the pre-level record and $\mathrm{wj}(D,j+1,\rho)$ in the
family. Keep the entire first term of (QW.17). For each pair
$(j,\rho)$ it contains $\log w_{j\rho}+\log(v/\pi)$.
There are exactly $f$ places and $m$ such entries per place.
The division by $f$ leaves $m\log(v/\pi)$ and at least $m$
from the middle weights. Adding $m(\log g_0+\log g_1)$ and
bounding $L+m^2L/k$ by $2L$ gives
$X_u\ge mA+m-2C_\epsilon L$. The exact lower endpoint of the
window subtracts $2\log\Gamma_u/(fk)$ on taking logarithms.
This proves (TFR.1). Under the additional bounds,
$\log T_u\ge mA$. Finally $2r_S=m-2$ and $A\ge0$ by $v\ge\pi$
and $g_i\ge1$. Exponentiating gives (TFR.2). Here $d_*(u)$ is
the stored size at the center of the window; the arithmetic column
dimension $D.\mathrm{duS}(u)$ does not enter this calculation.
\end{proof}
% END TFR1

% BEGIN TFR2
\paragraph{TFR2. Uniform frame and gap bounds along a supplied family.}
\label{claim:TFR2}
On every original family $\mathcal F$ with eventual
$\operatorname{vol}(Y_0)=\operatorname{massI}(D,D.R.r)$
at every index, (TFR.2) holds at every sufficiently late $x$,
for all indices $i\in\mathcal F.Ix(x)$ and words $u$ of
$D=D_{x,i}$. Independently of the volume identity, every supplied
family satisfies the original $\mathrm{GapHyp}(D_{x,i})$ at
every index beyond one family threshold.

\begin{proof}
Fix $0<\epsilon<1/52$ and let
$C_*=\max(5,\mathcal F.\mathrm{CYL}+1)$. DSZ4 verifies the
actual interlacing, endpoint and host bounds required by QW6 on
all words of the family. In particular
$4kx\Delta\le b_0<w_{0\rho}$ and $\Delta\ge1$.
The bottom coordinates increase and $b_j<w_{j\rho}$, so every
$w_{j\rho}\ge e$ for $x\ge1$ since $k\ge2$.
The rank window gives $m\ge x^{1/4}$, $L\ge1$ and $m^2L\le k$.
With $C_\epsilon$ as in (TFR.1), take
\[
 x\ge \max\{1,(2C_\epsilon+2)^{1/(1/4-\epsilon)}\}.
\]
Then $m\ge (2C_\epsilon+2)L\ge 2C_\epsilon L+2$.

For the mass terms, DSZ5 identifies the member's stored row
degrees with the exact Harish--Chandra polynomials in the
compact-dual measure. The supplied base-volume identity and the
family fields $\mathrm{FactDstar},\mathrm{FactDuStar}$ give
the two row and column expressions in (QW.14), with equality
before inserting the permitted row factor $(2/3)^f$.
The factorial, index and positive Euler product calculation of
QW1--QW6 therefore applies with these very $D_*,d_*(u)$.
The original $\mathrm{Gam\_eq}$ and rank-size bound, as in
DSZ9, give $0<\log\Gamma_u\le fk$. Take the maximum of the
displayed threshold, the volume threshold and those of DSZ4--DSZ5
and DSZ9. TFR1 now proves the assertion for the prime in the given
record at \emph{every} late index.

For the gap, write $n=D.R.r=k+m-1$ for the total row rank.
DSZ4 gives $k\ge3x$ eventually, so $n\ge1801$ eventually.
For every original place and adjacent bottom pair (DSZ.12) gives
\[
 b_{j+1}-b_j\ \ge\ x^{\epsilon/100} n^{19}f^5 L.
\]
Taking also $x\ge16^{100/\epsilon}$ ensures the right side is
at least $16n^{19}$. These are both clauses in the definition of
$\mathrm{GapHyp}$. The bounds used for this last assertion are
family rank and weight fields alone.
\end{proof}
% END TFR2

% BEGIN TFR3
\paragraph{TFR3. The closed frame event and the concrete selector.}
\label{claim:TFR3}
For any family class $\mathcal A$ with a supplied
$\mathrm{TSel}(\mathcal A)$, the original closed
$\mathrm{PX.TFrame}$ holds on $\mathcal A$, in particular at
\[
 (C_g,c_b,C_u,c_T,C_T,C_0)
   =(\mathcal F.\mathrm{CYL}+1,\tfrac14,3,\tfrac12,
                                    \mathcal F.\mathrm{CTu},3).
\]
Given also the compatible $\mathrm{SelS}$ of
$\mathrm{link\_alt}$, there is an original
$\mathrm{CSel}$ on $\mathcal A$, hence on $D_3(3,\mathcal A)$, consisting of concrete
packages at \emph{every} sufficiently late index. More generally a
supplied $\mathrm{SSelF}$ and the eventual base-volume identity
of TFR2 suffice for $\mathrm{CSelF}$ on a single family.

\begin{proof}
For a member of $\mathcal A$, let $(a,R)$ be its selected
$\mathrm{TSelF}$. For $x>\mathcal F.x_1$ and $x\ge a$ the field
$R_{x,i}.T.\theta.\mathrm{real}$ gives an original
$\mathrm{RealTied}$ for exactly $D_{x,i}$. RSL3 identifies its
$Y_0$ volume with the standard genus mass there. The threshold
$\max(a,\mathcal F.x_1+1)$ suffices for this input to TFR2.
The conclusion gives $\mathrm{TFrameAt}(D_{x,i})$ without an
admissibility premise. Thus it gives the implication required by
$\mathrm{PX.TFrame}$ for any six constants in its definition,
at the family's own $sC',sC_l$. An $x$ threshold supplied by
TFR2 works simultaneously for all $i$ and all words. In
$\mathrm{Closed}$ one may take $\epsilon_1=1$ and the zero
transported bound, since every family in the class already has
that eventual event, independently of its additional largeness test.
In particular restriction along $D_3(3,\mathcal A)\subseteq
\mathcal A$ gives precisely the frame premise of $\mathrm{link\_alt}$.

At any member of $\mathcal A$, the compatible $\mathrm{SelS}$
supplies spectral ties $(\mathrm{selS}_{x,i}).\mathrm{toSTies}$
starting at the \emph{same} $a$ as $R_{x,i}$. Take the maximum
of $a$, the rank-gap threshold and the frame threshold. Beyond
it TFR2 supplies both numerical inputs to SFS6 at every index.
Use its $\mathrm{ConcK}$ at the selected $X,Y$ for that very
index. The threshold and this dependent function of $x,i$ are
$\mathrm{CSelF}$. Applying the construction to every member gives
the asserted $\mathrm{CSel}$. For the single-family statement
use its supplied $\mathrm{SSelF}$ at its own threshold and the
given base-volume identity in the same argument. The
distinct-word estimate, diagonal terms and counts in the full
spectral comparison are additional inputs at their original scopes.
\end{proof}
% END TFR3

% BEGIN TFR4
\paragraph{TFR4. The frame row before a family is chosen.}
\label{claim:TFR4}
Fix the source parameters $h_L:\mathrm{LocalLanglandsGL}$ and
$p:\mathrm{Bl.Prm}(h_L)$. Then the original row assertion
\[
\mathrm{RowE}(\mathrm{Bl.IsGoodE}(p.s),\mathrm{Bl.IsLinked}(p),
 \mathrm{TFrameAt})
\]
holds. In fact the estimate uses only the
construction context and the fact that the point lies over its
pre-level data. This is the row quantified over all compatible
points and bundles \emph{after} one common $x$ threshold.

\begin{proof}
Fix $0<\epsilon<1/52$, a standard constant record $C$, and the
construction exponents $c_D,c_h\ge0$. Let $K$ be any context
with these values at scale $x$. It supplies $P=K.P$, the level
choice $K.\mathrm{Lv}$ and $K.\mathrm{ok}:\mathrm{ChOK}$ for its
original coordinates and words. In (TFR.1) use $C_*=5$, hence
$C_\epsilon=30+12/\epsilon$. QW7--QW9 verify the place, rank
and discriminant hypotheses uniformly before $K$ or a word is
given. In particular $k\le x^5$, $m^2\le k$, $v\ge\pi$,
$w_{j\rho}>b_0\ge4kx\Delta$ and $m\ge x^{1/4}$ eventually.

The two size equations in $K$ are precisely
\begin{align*}
 d_*(u)&=\mu_k I_{\rm col}\prod_\rho d_k(u_\rho),\\
 D_*&=\mu_{k+m-1} I_{\rm row}\prod_\rho
                  d_{k+m-1}(v_\rho,-b_\rho),
\end{align*}
where $d_t(z)=\prod_{i<j}(z_i-z_j)/\prod_{i=0}^{t-1}i!$.
For the second equation the source defines
$\mathrm{trueDsig}$ by the formal-degree \emph{polynomial}.
Its entries are exactly the appended $v,-b$ in decreasing order,
and collecting denominator factors by their larger index gives
$\prod_{j=0}^{k+m-2}j!$. The pre-level record's $v,b$ agree with
those of the base word. Thus the stated equation holds for every
word and implies the required row lower bound with $(2/3)^f$.
The two level indices here are exactly those counted in QW2,
with $f$ copies of each of the two split-prime factors. The
standard masses are the ones in MF19 and the telescope of QW3.
This verifies the arithmetic inputs to (TFR.1) on $K$ itself.

There is one difference in the window comparison on a raw
context. Its actual $\Gamma$ is
$\exp(C_\Gamma m^5 fL)$, and $K.\mathrm{Lv}.\mathrm{T\_win}$
applies to this value. The rank fields give
$m\le2\sqrt{x}$ because $\delta_1=x^{-\epsilon/8}\le1$,
and $k\ge x^{9/2+10\epsilon}\ge x^{9/2}$. Hence
\[
 \frac{2\log\Gamma}{fk}
       \le 64|C_\Gamma|\,x^{-2+\epsilon}.
\]
This is at most $2$ once
$x\ge\max(1,(32|C_\Gamma|)^{1/(2-\epsilon)})$.
Take also the threshold of TFR2's elementary comparison
$x^{1/4}\ge2C_\epsilon L+2$, and those of QW7--QW9.
TFR1 now gives (TFR.2) for each prime $K.\mathrm{Lv}.T(u)$.

For an actual point $\beta$ over $K$, its core field
$\beta.c.\mathrm{isOver}$ identifies the stage base with
$\mathrm{Pre.ofCh}(P,K.\mathrm{Lv},\ldots)$. This construction
takes its words and ranks from $P$ and its prime function from
$K.\mathrm{Lv}$. The enlarged level-prime set in
$K.\mathrm{Lv}.\mathrm{host}$ preserves $g_0,g_1$ and the
underlying field defining $v$. The stage's column field
$\mathrm{A\_T}(u)$ identifies the member's $T_u$ with the base
prime at $u$. Finally $\beta.D$ is obtained from this stage by
$\mathrm{setDiv}$, which updates the division lattices and
their heights. Projecting the words, $T_u$, ranks and host from
its definition gives the quantities just bounded. Thus (TFR.2)
is the demanded $\mathrm{TFrameAt}(\beta.D)$.

For the outer $\mathrm{RowE}$ quantifiers one can use
$\epsilon_0=1$, the zero transported bound, and lower thresholds
$c_1=c_2=0$ for $c_D,c_h$. After the standard $C$, $c_D$ and
$c_h$ have been given, choose the finite maximum of the thresholds
above. The comparison then holds on every $K$ with
$\kappa.\mathrm{Of}(K)$ at that scale, every eligible point,
each real $c$ and every linked bundle over it. This gives the
original row for the supplied $h_L,p$. On this construction
route the row uses the written polynomial $\mathrm{trueDsig}$.
On general families TFR2 instead uses DSZ5 to identify their
stored analytic degrees.
\end{proof}
% END TFR4

\section{The absolutely convergent kernel on the selected family}
Fix a datum $D$ and a spectral tie $T$ on it. This records a setting
$X$, full arithmetic data $Y$, and, for every proof $h_F$ of the
finiteness record of $Y$, a comparison from the original spectral
datum $D.s$ to the constructed datum $s_Y=\mathrm{sOf}(h_F)$.
The expansion below is an identity on the original column spaces
of $D.s$, including its normalizing scalar.

% BEGIN KXE1
\paragraph{KXE1. The dimension in the normalizing constant.}
For any supplied setting $X$ and word $u$, the dimension of its
compact representation $V_u$ is the product $\mathrm{dimPi}(u)$
in the definition of the abstract Gram kernel.
\begin{proof}
At a place $v$ let $\lambda$ be the compact Harish--Chandra parameter
attached by $X.\mathrm{pl}$ to $u$. Its entries are strictly decreasing.
The compact Weyl dimension formula gives, in these shifted coordinates,
\[
            \dim V_{u,v}=
        \frac{\prod_{a<b}(\lambda_a-\lambda_b)}
             {\prod_{a=0}^{k-1}a!}.
\]
Indeed the positive roots are the differences of coordinates,
and their pairings with the half-sum of positive roots give
$\prod_{a<b}(b-a)=\prod_{a=0}^{k-1}a!$. The source defines its
local $\mathrm{dimPiU}$ as the least dimension of a representation
of this type. The given $V_{u,v}$ belongs to the defining set and
the same formula gives the same dimension for every member.
This proves the local identity.

Choose orthonormal bases $b_{v,j}$ and take the supplied tensor
vectors $e_{\mathbf j}=\mathrm{tens}_u((b_{v,j_v})_v)$.
The product inner-product axiom makes them orthonormal. For
$y=\mathrm{tens}_u((x_v)_v)$, local Parseval and the same axiom give
\[
 \sum_{\mathbf j}|\langle e_{\mathbf j},y\rangle|^2
 =\prod_v\sum_j|\langle b_{v,j},x_v\rangle|^2
 =\prod_v\|x_v\|^2=\|y\|^2 .
\]
Consequently $y$ belongs to their span, by orthogonal projection.
All the supplied tensors span $V_u$, so these particular ones form
a basis. Counting the multi-indices and using the given bijection
from places to the indices of $\mathrm{dimPi}$ gives the assertion.
This proof uses the inner-product and spanning fields even when the
tensor map is presented without multilinearity axioms.
\end{proof}
% END KXE1

% BEGIN KXE2
\paragraph{KXE2. The original kernel equals the sum of its terms.}
Assume the supplied $X=T.X$ satisfies $\mathrm{WuDecay}$ and its
archimedean kernel is integrable. For all ordered word pairs $(u,u')$,
\[
 K_{D.s}(u,u')=\sum_t A_t,\qquad
                  \sum_t\|A_t\|<\infty ,
 \qquad A_t=(D.s).\mathrm{aTerm}(u,u',t).
\]
The operator sum converges as a net over finite subsets. This is
$\mathrm{PX.Exp}(D)$. The source gap on $D$ suffices for both
analytic assumptions.
\begin{proof}
SFS2 gives the level normalization for these $X,Y$. In the proof of
QPR11 the three hypotheses extracted from the concrete tie were
that normalization, the decay of the supplied wave and integrability
of $X.\mathrm{ph}$. To see their sufficiency on the given data,
FLD3 folds both word columns, and QPR3 places their period series
in the same Hilbert range. QPR4 chooses strict domains, QPR5--QPR8
give the orbital integrals and the absolutely summable norm bounds,
and QPR9--QPR10 give independence of the domains and column
representatives. PGR18 then gives the ungrouped operator series with
its normalized Gram pairing for $Y$. This argument works for equal
words as well. SFS5 now supplies $h_F:\mathrm{FinX}(Y)$.
Use KXE1 at each of the two words and apply PGR20. Its finite-fibre
regrouping gives the asserted sum and absolute convergence first
on $s_Y$.

Take $I=T.hI(h_F)$, writing $U_w=I.eCu(w)$ and $U_H=I.eHH$
for its unitary identifications. Conjugation
$A\mapsto U_u A U_{u'}^{-1}$ preserves operator norm and finite
sums. The term comparison in $I$ sends each $A_t$ to the term on
$s_Y$ indexed by $I.eT(t)$, a bijection for this pair of words.
The period-map comparison is $U_HP_w=P'_w U_w$. Taking adjoints
gives
\[
 U_u P_u^*P_{u'}U_{u'}^{-1}=(P'_u)^*P'_{u'}.
\]
For $s=D.s$ or $s=s_Y$, write
\[
 c_w=s.\mathrm{volN}(w)s.\mathrm{volH}(w)/\mathrm{dimPi}(w).
\]
The volumes agree across $I$ and the parameters in $\mathrm{dimPi}$
are shared. Thus $(c_uc_{u'})^{-1/2}$, the real scalar multiplying
the operator product, agrees as well. The original $K_{D.s}$ maps
to $K_{s_Y}$. Pull the convergent finite sums back by inverse
conjugation and reindex their finite subsets by the inverse bijection;
the norm series pulls back with the same terms. This proves both
assertions on $D.s$. At a datum with the original gap SFS4 gives
both analytic hypotheses on $T.X$ itself.
\end{proof}
% END KXE2

% BEGIN KXE3
\paragraph{KXE3. One threshold for the kernel expansion at every index.}
An original family with a supplied $\mathrm{SSelF}$ has a single
threshold above which $\mathrm{PX.Exp}(D_{x,i})$ holds for
\emph{every} index $i$. In particular a class with a supplied
$\mathrm{SSel}$ satisfies the original closed $\mathrm{PX.ExpF}$.
\begin{proof}
Let $a$ be the threshold of the supplied selector, $b$ the threshold
for the gap in TFR2, and $x_1$ the strict family cutoff. At each
$x\ge\max(a,b,x_1+1)$ and each $i$, use the original $T_{x,i}$
and its own comparison. TFR2 supplies the gap at this index;
KXE2 supplies $\mathrm{PX.Exp}(D_{x,i})$ for every pair of words.
This choice of threshold works simultaneously for all the parameters
$C_0,C_g,c_b,C_u,c_T,C_T$ in $\mathrm{PX.Under}$. In particular
it works when its standing hypotheses hold with the family's own
$C',C_l$, giving precisely $\mathrm{ExpF}$. The argument on a class
takes the selector at the given family and its class-membership proof.
It works for every family of the class, so any of the outer closed
size thresholds for the identically true predicate suffice here.
\end{proof}
% END KXE3

\section{Comparing the strict Selmer lengths on the chosen frame}
\subsection{Passing from the frame's finite coefficients to the arithmetic length}
\label{sec:host-selmer-length}

% BEGIN VLT1
\paragraph{VLT1. The strict Selmer group on the chosen quotient.}
\label{claim:VLT1}
Let $\Omega$ be the host's algebraic closure of $\mathbb Q$, $\ell$ its
prime, $O=c.O$ its coefficient DVR, $N>0$ its level, $e$ its exponent,
and $\psi$ its Dirichlet character. Assume the source hypotheses
$\mathrm{HostData}(\ell,e,\psi,o_1)$ and
$\mathrm{IsHostFrame}(\Omega,c,e,\psi,S,A,F)$.
Write $K=\operatorname{Frac}(O)$, $M=K/O$ with the discrete topology,
$G=\operatorname{Gal}(\Omega/\mathbb Q)$ and $\Gamma=G_S$.
The set $S$ consists of the primes dividing $\ell N$. The frame's map
$\pi:G\to\Gamma$ is the quotient by the closed normal subgroup generated
by the inertia groups above primes outside $S$. For each $v\in S$ write
$D_v,I_v\subseteq G$ for the decomposition and inertia groups of the
chosen $A_v$, and $\mathrm{loc}_v=\pi|_{D_v}$. The action on $M$ is by
$\alpha=\Psi^{-1}$ on $\Gamma$, where
$\Psi\circ\pi=\epsilon_\ell^{\,e}\psi_{\rm Gal}$.
The action is continuous since the cyclotomic character is continuous,
the Dirichlet character has open kernel, and every element of $M$ is
killed by a power of the maximal ideal. RV1 identifies $M$ with the
source's $O\otimes_{\mathbb Z_\ell}(\mathbb Q_\ell/\mathbb Z_\ell)$.

Define
\[
  \mathcal S=\bigcap_{v\in S}
   \ker\!\left(H^1_{\rm cts}(\Gamma,M(\alpha))\longrightarrow
   H^1_{\rm cts}(D_v,M(\alpha\circ\mathrm{loc}_v))\right).
\]
Then $\operatorname{length}_O\mathcal S=\mathrm{Vlen}(\Omega,c,e,\psi)$.
Indeed, the source's absolute local conditions for $\mathrm{Vlen}$ are
vanishing on $I_B$ for every valuation ring $B$ of $\Omega$, and on $D_B$
when $B(\ell N)<1$. The latter rings lie over primes of $S$.
An absolute cocycle satisfying these conditions vanishes pointwise on
every inertia group outside $S$, since the character is trivial there
and every local coboundary is zero. The action on $\ker\pi$ is trivial,
so the restriction of the cocycle to this kernel is a continuous
homomorphism to the additive group of $M$. It vanishes on the generators
and their conjugates (the conjugates are again such inertia groups),
hence on the closed subgroup they generate. The cocycle therefore
descends through $\pi$; continuity follows from the quotient topology.
The same formula $(\alpha(g)-1)m$ describes global coboundaries on both
sides after inflation, and surjectivity of $\pi$ gives injectivity on
classes. The descended class restricts to zero on every selected $D_v$.

Conversely, local coboundaries on the selected groups remain
coboundaries on their conjugates. In fact if an inflated cocycle
$z(d)=(\alpha(\pi d)-1)m$ on $D_v$, then on $gD_vg^{-1}$ its primitive
is $\alpha(\pi g)m-z(g)$ by the cocycle identity. All valuation rings
over the same rational prime are conjugate in the algebraic closure.
Thus inflation gives the required condition at every $D_B$ over $S$,
and hence at its inertia. Inertia outside $S$ maps trivially under
$\pi$. A valuation on an algebraic extension of $\mathbb Q$ that is
trivial on $\mathbb Q$ is trivial everywhere, with trivial inertia and
$B(\ell N)=1$. This exhausts the absolute conditions and gives the
claimed $O$-linear isomorphism and length identity.
% END VLT1

% BEGIN VLT2
\paragraph{VLT2. Lifting a torsion class with the frame's local conditions.}
\label{claim:VLT2}
Put $T_b=\{v_\ell\}\cup F.S_b$. At a prime $p\ne\ell$ of $S$,
membership in $F.S_b$ means that the $p$-component of the
$\ell$-power part $\kappa$ of $\psi$ is nontrivial. Put
$N_v=I_v\cap\ker(\eta\circ\pi)\cap\ker(\bar\Psi\circ\pi)$ inside $D_v$.
We first record the two local consequences needed below.

If $v\notin T_b$, the cyclotomic character is unramified there.
The character $\kappa$ is also unramified since it factors away from
the $p$-part of $N$ at this place: inertia fixes all roots of unity
of order prime to $p$. Consequently $\Psi(\pi s)$ on inertia has
order prime to $\ell$. For $s\in N_v$ its residue equals one. If a
unit $u$ of prime-to-$\ell$ order $b$ has residue one, the sum
$1+u+\cdots+u^{b-1}$ reduces to the nonzero residue of $b$, so
$(u-1)(1+u+\cdots+u^{b-1})=0$ implies $u=1$. Thus $\Psi(\pi N_v)=1$.

If $v\in T_b$, there is $x\in D_v$ for which $\alpha(\pi x)-1$
is a unit of $O$. At $v_\ell$ the element of HSR3 has
$\Psi(\pi x)^2-1$ a unit. This makes $\Psi(\pi x)-1$ and hence
$\alpha(\pi x)-1$ units. For another $v\in T_b$ write $p$ for its
prime. A nontrivial $\ell$-primary character of
$(\mathbb Z/p^{v_p(N)})^\times$ implies $\ell\mid p-1$ since
$p\ne\ell$. Thus $p$ is odd. Choose an arithmetic Frobenius lift
$\sigma\in D_v$; it acts on all prime-to-$p$ roots of unity by
raising to $p$. Hence $\bar\epsilon_\ell(\sigma)=1$ and the unit
$u$ describing its action on $\mu_N$ has dyadic coordinate $p$.

Here is the dyadic calculation at this $p$. Write $t=c_2(\psi)$,
$o=o_1\ge1$, $s=v_2(p^2-1)\ge3$ and $j=v_2(o)$.
Squaring an integer $1+2^s b$ with $b$ odd raises the exact valuation
of its difference from one by one, and an odd power preserves that
valuation. Thus $d=v_2(p^{4o}-1)=s+j+1$. The host gives
$2^d\le 2o(p^2-1)<8op^2<2^t$.
The element $p^{4o}$ generates $1+2^d\mathbb Z/2^a\mathbb Z$,
where $a=v_2(N)\ge t$, by the same squaring formula for its order.
If $\psi_2(u)^{4o}=1$, the dyadic component would factor modulo
$2^d$, contradicting its conductor $2^t$. The host requires that each
odd-prime component of the prime-to-$\ell$ part $\theta_0$ of $\psi$
is killed by $2o$. HXD1 identifies its dyadic component with $\psi_2$.
It follows that $\theta_0(u)^{4o}=\psi_2(u)^{4o}\ne1$.
Reduction kills $\kappa(u)$ and is injective on prime-to-$\ell$ roots
as shown above. Thus $\bar\psi(u)^4\ne1$. Taking $x=\sigma^2$ gives
$\bar\Psi(\pi x)^2=\bar\psi(u)^4\ne1$, which proves the asserted unit
condition at this place as well.

For $0\ne a\in O$, write $L_a=O/(a)$ with action $\alpha$, and let
$j_a:L_a\hookrightarrow M$ carry $r$ to $r/a+O$; its image is $M[a]$.
The source's $F.\mathrm{ShaO}(a)$ is the $O$-span in
$H^1_{\rm cts}(\Gamma,L_a)$ of classes having a cocycle representative
zero on $\pi N_v$ for $v\notin T_b$ and a coboundary on $D_v$ for
$v\in T_b$. Its extended $O$-length is $F.\mathrm{shaLen}(a)$.
We have
\[
    \mathcal S[a]\ \subseteq\ (j_a)_*\big(F.\mathrm{ShaO}(a)\big)
    \quad\text{inside }H^1_{\rm cts}(\Gamma,M(\alpha)).
\]
To see this, choose a cocycle $z$ representing a class in $\mathcal S[a]$.
Write $az=\delta m_0$, with $\delta m(g)=(\alpha(g)-1)m$.
Divisibility of $K/O$ gives $m_0=am_1$. The cohomologous cocycle
$z-\delta m_1$ takes values in $M[a]$ and therefore uniquely has the
form $j_a y$. The map $y$ is a continuous cocycle since $j_a$ is
injective, equivariant and both coefficient modules are discrete.
For each $v$ choose a primitive of $z|_{D_v}$ and subtract $m_1$ from
it, obtaining $b_v\in M$ with
$j_a y(\pi d)=(\alpha(\pi d)-1)b_v$ for all $d\in D_v$.
Outside $T_b$ this gives $y(\pi N_v)=0$ by the first local calculation.
Inside $T_b$ evaluate at the element $x$ of the second calculation.
Multiplication by $a$ gives $(\alpha(\pi x)-1)ab_v=0$.
Since the first factor is a unit, $ab_v=0$. Hence $b_v=j_a b_v'$
with $b_v'\in L_a$, and injectivity gives $y|_{D_v}=\delta b_v'$
through $\mathrm{loc}_v$. This places $[y]$ in the defining span and
$(j_a)_*[y]=[z]$. Inclusion into the image, followed by comparison
with its source, proves
$\operatorname{length}_O\mathcal S[a]\le F.\mathrm{shaLen}(a)$.
% END VLT2

% BEGIN VLT3
\paragraph{VLT3. The literal host length comparison.}
\label{claim:VLT3}
For every natural number $V$ satisfying
$F.\mathrm{shaLen}(a)\le V$ for all $0\ne a\in O$, one has
\[
  \mathrm{Vlen}(\Omega,c,e,\psi)\le V.
\]
Indeed $\Gamma$ is a quotient of the compact group $G$ and the image
of each continuous cocycle $\Gamma\to M$ is finite. Each value is
annihilated by a power of $\ell$ since $v_O(\ell)>0$. Taking the
largest of finitely many required exponents kills the cocycle and its
class. Thus the increasing submodules $P_k=\mathcal S[\ell^k]$ have
union $\mathcal S$. VLT2 applies at every nonzero $\ell^k$, including
$\ell^0=1$, and gives $\operatorname{length}_O P_k\le V$. Choose an
index attaining the largest length among these integers. Every later
inclusion is an equality because its finite-length quotient has
length zero. The chosen $P_k$ therefore equals $\mathcal S$. VLT1
now gives the claimed bound, including $V=0$.

In particular $\mathrm{Vlen}$ is finite before the operation
$\mathrm{toNat}$ in the source's real normalization.
With $e_O=\operatorname{length}_O(O/\ell O)=v_O(\ell)>0$ from RV1,
the contribution to $\mathrm{VLe}$ consequently satisfies
\[
   \mathrm{Vn}(\Omega,c,e,\psi)\log\ell
      =\frac{\operatorname{length}_O\mathcal S}{e_O}\log\ell
      \le \frac{V}{e_O}\log\ell .
\]
Here $V$ is precisely the uniform upper bound on the host's finite
coefficient lengths. It need not be attained by one of those lengths.
% END VLT3

\section{From the Case V guard to a separating trace bound}
\label{sec:count-v-install-v61}

Here $\ell$ is the original coefficient prime, $O=c.O$, $k$ its residue
field, and $K=\mathbb Q(i)$. The source supplies a specified algebraic closure $\Omega$ of $\mathbb Q$. Write
$\Delta=\operatorname{Gal}(\Omega/\mathbb Q)$.
The algebraic inputs called Lemma M and Lemma M-prime in this section
are the valuation and length results SP1--SP6.

% BEGIN CVI1
\paragraph{CVI1. The Case V hypotheses supply both host predicates.}
\label{claim:CVI1}
Fix a character $\chi$, a host $h$ with $h.\chi=\chi$, a permitted
$t$ with $h.\mathrm{Big}(t)$, an odd prime $\ell$, and
$1\le e<t$. Fix $\psi\in h.X(c,e)$ and suppose that
$h.\mathrm{CaseV}(\ell,e)$ and
$R_1=h.R1(\ell,\psi)\le2^{c_2(\psi)}$.
In particular $e=1$ and $\ell\in B_\chi$. Take $q=h.q$, $o=h.o$,
$N=h.N$; $q$ is the positive conductor of $\chi$, and $o$ its order.
Put $\theta=\operatorname{th1}(\ell,\psi)$ as in HXD1.

Here is the character calculation at this host. Its defining subgroup
is the intersection of the inverse images of $\langle-1\rangle$ at
the dyadic level, $\ker\chi$ at $q$, and the square subgroups at the
two auxiliary odd primes. It is killed by $\psi$. For each odd
$p\mid q$ the full $p$-power in $N$ is the one in $q$, since the
two auxiliary primes are coprime to $q$. The $o$-th power of every
unit supported at $p$ lies in the host subgroup. At each of the two
auxiliary primes outside $q$ the square of every local unit lies in
that subgroup. At the other odd primes outside $q$ there is no
component. Thus the odd $p$-parts of $\psi$, and their prime-to-$\ell$
parts, have orders dividing $2o$. The square of $\theta$ has conductor
dividing $q$: its dyadic component is trivial by construction,
its components at the two auxiliary primes square to one, and the
remaining conductor powers divide $q$. The $\ell$-primary part of
$\psi$ is trivial on the dyadic and auxiliary factors since those
images have order prime to $\ell$. In particular every
$v\in T_b(\ell,\psi)$ divides $\ell q$.

The order calculation of HXD1 uses just the odd prime and the prime-to-$\ell$ order of $\theta$: taking $d$ to be the order of $\ell$ modulo $\operatorname{lcm}(4,\operatorname{ord}\theta)$ gives $n_0=v_2(\ell^d-1)\ge2$. Since $o,q\ge1$ and
$\ell q\ge3$, the original $R_1=16o(\ell q)^2 4^{n_0}$ is at least
$16\cdot9\cdot16>2^{11}$. The guard implies $c_2(\psi)\ge12$
and for every $v\in T_b$,
\[
 8o v^2\le8o(\ell q)^2<R_1\le2^{c_2(\psi)} .
\]
Moreover $16\mid N$ by the original host exponent $\mathrm{ex}\ge2$,
and $\psi(-1)=(-1)^e$ by membership in $h.X(c,e)$. These are the
fields of $\mathrm{HostData}(\ell,e,\psi,o)$, including the
prime-to-$\ell$ component bound above. Since
$2(\ell q)4^{n_0}\le R_1$, the positive integer $q$ and the
conductor divisibility give $\mathrm{XData}(\ell,\psi,q)$.
All calculations allow $\ell=3$ and equality at the $R_1$ boundary.
% END CVI1

% BEGIN CVI2
\paragraph{CVI2. One actual Galois frame and the class-one condition.}
\label{claim:CVI2}
Let $S$ consist of the rational primes dividing $\ell N$, choose
one valuation ring $A_v$ above each, and let $\Gamma=G_S$ be the
quotient of $\Delta$ by the closed normal subgroup generated by all
inertias above primes outside $S$. Use its quotient map $\pi$ and
the actual $D_v=D(A_v)$ and $I_v=I(A_v)$. Write $v_\ell$ for the
place at $\ell$ and use the supplied $S_b(\ell,\psi)$ for the marked
set of bad places. The Galois character $\eta_4$ of $\mathbb Q(i)$
and $\chi_\ell^e\psi_G$ descend to $\Gamma$: at every omitted
prime their finite-level cyclotomic factors are unramified, and their
kernels are closed. Call the descended characters $\eta,\Psi$.
The image of a complex conjugation is the element $c_F$ of the frame.
This gives every field of $\mathrm{IsHostFrame}$. The order of choices
keeps $S,A_v,F$ fixed for the rest of the argument.

The native $F.\mathrm{Ok}$ conditions hold. The residue field is
finite of characteristic $\ell$ by HXD1, and $2$ is a unit. The
quotient map is continuous, and the finite-order character and
all reductions of $\Psi$ modulo powers of the maximal ideal are
locally constant. On $c_F$ their respective values are $-1,1$ by
the parity in CVI1, and $c_F^2=1$. The inertia and decomposition
groups are closed, with inertia normal in decomposition. Their
quotient is generated topologically by residue Frobenius and maps
onto every finite cyclic group with Frobenius mapping to $1$:
the reductions on the finite residue extensions give exactly these
maps. At $T_b=\{v_\ell\}\cup S_b$ the primes are odd (the
dyadic unit group has $2$-power order), so $\eta$ is unramified.
AFE10 provides, in each of these actual $D_v$, an $x$ with
$\eta(x)=1$, $\bar\chi_\ell(x)=1$ and $\bar\Psi(x)^2\ne1$.
Here and below characters on $\Gamma$ are pulled back along $\pi$
when evaluated on $D_v$. Away from $T_b$ the cyclotomic factor
is unramified and the ramified finite part of $\psi$ has order
prime to $\ell$. Prime-to-$\ell$ roots of unity reduce injectively.
Consequently $\Psi=1$ on the prescribed joint inertia kernel
$N_v=I_v\cap\ker\eta\cap\ker\bar\Psi$. Finally $v_\ell\notin S_b$
by definition. These verifications include all sixteen frame conditions.

Both $\Delta$ and $\Gamma$ are profinite. Global finiteness implies
$\Phi_\ell(\Gamma)$. Indeed, for any open $H\le\Gamma$, induce the
trivial $\mathbb F_\ell$-module of $H$ to the finite permutation
module on $\Gamma/H$. Shapiro identifies its continuous $H^1$ with
$\operatorname{Hom}_{\rm cont}(H,\mathbb F_\ell)$; finite coefficients
and $S\supseteq\{\ell\}$ give finiteness of the former group.

For $F.\mathrm{Cl1}$ let $a:\ker\eta\to B$ be any locally constant
homomorphism to an abelian group which kills the specified inertia
conjugates. Pull back to $G_K\subset\Delta$, the inverse image of
$\ker\eta$. The homomorphism has finite image by compactness.
Every inertia of $K$ is an intersection of $G_K$ with an absolute
inertia of $\mathbb Q$. Outside $S$ the latter group dies in
$\Gamma$; inside $S$ it is conjugate to the chosen $I_v$, and the
conjugate kernel condition on $a$ kills the intersection. Thus the
finite abelian extension defined by the pullback is everywhere
unramified. The Gaussian integers are Euclidean; by the unramified
abelian class-field theorem this extension is trivial. Since
$G_K$ surjects onto $\ker\eta$, $a$ is trivial.
% END CVI2

% BEGIN CVI3
\paragraph{CVI3. Local vanishing and the global degree-two restriction.}
\label{claim:CVI3}
For every integer $j$ put $M_j=\mathrm{Gb}(k,j)$, the matrices
supported at entries of weight $j$ in the three-dimensional
residual adjoint action. The dimension is
$m_j=\max(3-|j|,0)$. For $j\ne0$ the AFE10 witness of CVI2 acts
on $M_j$ by $\bar\Psi(x)^j$, so $x-1$ acts invertibly:
only $j=\pm1,\pm2$ can contribute. It acts invertibly on
$M_j^*(1)$ as well, since $\bar\chi_\ell(x)=1$.
Thus at $v\in T_b$, the degree-zero local cohomology of the
module and of its Tate dual vanishes. Local duality gives
$H^2(D_v,M_j)=0$. At $v\in S_b$, the local Euler characteristic
away from $\ell$ then gives $H^1(D_v,M_j)=0$. These are exactly
the two clauses of $F.\mathrm{LocVan}$.

For $j\ne0$ the same argument at $v_\ell$ kills the global
Tate-dual invariants. For $j=0$, choose a Galois automorphism
fixing $i$ and acting nontrivially on $\mu_\ell$. It exists by
the CRT at $4$ and $\ell$ and surjectivity of the cyclotomic
quotient over $\mathbb Q$. On degree-zero matrices its action
is the identity (its $\eta$ is $1$), whereas its action on
$M_0^*(1)$ is a nonidentity scalar. Poitou--Tate therefore makes
\[
             H^2(\Gamma,M_j)\longrightarrow
                     \bigoplus_{v\in S} H^2(D_v,M_j)
\]
surjective for every $j$. Odd residue characteristic makes the
positive-degree archimedean Tate cohomology zero. The local and
global finiteness assumptions apply to the original finite
$\mathbb F_\ell$-modules by restriction of scalars from $k$;
the underlying cochain complexes and their maps are unchanged.
Composition with $\operatorname{Tr}_{k/\mathbb F_\ell}$ identifies the
restricted $k$-linear dual with the $\mathbb F_\ell$-linear dual by
nondegeneracy of the finite-field trace pairing, compatibly with
the cyclotomic twist.
% END CVI3

% BEGIN CVI4
\paragraph{CVI4. The Euler defect for the native framed quotient.}
\label{claim:CVI4}
Write $Z_{L,v}\subseteq Z^1_v$ for the local subspace at weight $j$.
All local coboundaries lie in it: off $v_\ell$ the action on $N_v$
is trivial, and at $v_\ell$ the only additional weight-zero
inertia condition is trivial for a coboundary. Put
$L_v=Z_{L,v}/B^1_v$, $U=\bigoplus_vH^1_v/L_v$.
Localization gives an exact sequence
\[
 0\to H^1_L\to H^1\to U\to
          F.\mathrm{Ob}(j)\to H^2\to\bigoplus_v H^2_v\to0.
\]
Indeed an obstruction representative is a global $2$-cocycle and
a local $1$-cochain at each $v$ whose differential equals its
restriction. Forgetting the local cochains maps onto the kernel of
the last restriction. When the global cocycle is a coboundary,
subtract a global primitive together with its restrictions. What
remains is a tuple of local cocycles modulo $Z_{L,v}$ and global
restrictions, precisely the indicated cokernel. CVI3 supplies
surjectivity at the last term. The numerator defining $F$'s
degree-one quotient is the space of cocycles with precisely these
local conditions, giving $H^1_L$ at the other end. Finite-dimensionality
follows at once. Its Euler difference is
\[
 h^1-h^2+\sum_v(h^2_v-\dim Z^1_v+\dim Z_{L,v}).
\]

For $v\ne v_\ell$ the quotient $I_v/N_v$ has order prime to
$\ell$ by the joint-character argument in HXP1, and $N_v$ acts
trivially on $M_j$. Locally $L_v$ is the kernel of restriction to
$N_v$, equivalently to $I_v$: on $N_v$ coboundaries are zero;
if a class vanishes there, its restriction to $I_v$ inflates from a
finite group of invertible order and vanishes by averaging. The
unramified local space has dimension $h^0_v$, since the unramified
quotient is procyclic and its first cohomology on $M_j^{I_v}$
is the Frobenius cokernel, of the same dimension as the kernel.
It follows that $\dim Z_{L,v}=m_j$, cancelling the local term by
the Euler formula $h^1_v=h^0_v+h^2_v$.

Let $a_j=\dim M_j^{c_F}$. Negative transpose followed by reversal
pairs the entries $(a,b)$ and $(2-b,2-a)$; its fixed space has
dimension $a_j=1$ for $|j|\le1$ and $0$ otherwise, since the
characteristic is odd. The global Euler formula gives
$h^1-h^2=h^0+m_j-a_j$ and the $\ell$-adic local formula gives
$h^1_\ell=h^0_\ell+h^2_\ell+m_j$. For $j<0$ the local
condition is $B^1_\ell$; both global and distinguished local
invariants vanish by CVI3, so the difference is $-a_j$. For
$j>0$ the local condition is all of $Z^1_\ell$, $h^2_\ell=0$,
and the difference is $m_j-a_j$. For $j=0$ it is unramified,
with $\dim Z_{L,\ell}=m_0$ by the same Frobenius computation
($I_\ell$ acts trivially). The global invariants at weight zero
are exactly $M_0^{c_F}$ since the sole action on the diagonal
comes from $\eta$. The difference is zero. Thus the integer
$\dim H^1_L-\dim F.\mathrm{Ob}(j)$ equals $-1$ at $j=-1$,
$1$ at $j=1,2$, and zero otherwise. This is the original $F.\mathrm{Count}$.
% END CVI4

% BEGIN CVI5
\paragraph{CVI5. The finite list obtained from a compatible weighted model.}
\label{claim:CVI5}
Fix the original $\iota:\overline{\mathbb Q}_\ell\simeq\mathbb C$.
From $\mathrm{ThmS6\_s}$ take its $g_1\in\Gamma$, marked on
the frame of CVI2. Suppose the source's native marked universal
pair $(R,u)$ and its weighted deformation comparisons described below
are supplied. Then the two algebraic lemmas give an integer $V$ and
a finite list $r_i$ of marked characteristic-zero points of type I
satisfying $B'_s(K,\iota)$ on $\Delta$ and pairwise nonisomorphic
on $\ker\eta$, with
\[
  \mathrm{shaLen}_F(b)\le V\quad(0\ne b\in O),\qquad
  V\le2 e_O\sum_i v_\ell\bigl(
  \operatorname{tr}r_i(\phi)-\operatorname{tr}r_z(\phi)\bigr).
\]
Here $e_O=v_O(\ell)$, $v_\ell(\ell)=1$,
$r_z$ denotes the original diagonal base representation, and
$\phi$ is the single separating element supplied by
$\mathrm{HostNumber\_s}$ for $g_1$. We specify the comparisons used
in this implication to identify the exact scope of the model assumption.

Take an adapted minimal weighted power-series presentation $W\to R$
with origin $u\mapsto r_z$. It must identify its weight parts with
the native parts of the universal representation. The variables
are marked matrix-entry differences from the base representation,
of weights in $[-2,2]$. Write $d_j$ for their
numbers. The native obstruction-to-relations and tangent-to-generators
comparisons give homogeneous kernel generators, at most
$b_j=\dim F.\mathrm{Ob}(j)$ in weight $j$, and
$\dim H^1_L(j)\le d_j$. CVI4 then gives
$b_j+\mathrm{defect}(j)\le d_j$. The nonpositive tangent
comparison takes $F.\mathrm{VanE}$ to the assertion that an
$O$-valued coordinate vector of any weight $j\le0$ killing the
linear parts of these relations is zero. These are comparisons
of the actual framed representations with the chosen coordinates.

The needed point comparisons are as follows. On the graded
entry ring $\mathrm{Sring}(u)$ each geometric point $\xi$ gives
the representation $r_\xi$ by evaluating all entries, with its
original right component. Finiteness of the weight-zero module
gives a marked characteristic-zero point after permitted torus
rescaling. Vanishing on all negative or all positive weights
corresponds to the two triangular flags. Vanishing on both gives
the origin, by the native diagonal-point comparison using $F.\mathrm{Cl1}$.
A point having a nonzero coordinate on each side gives type I,
by the native two-sign comparison using $F.\mathrm{LocVan}$
and $F.\mathrm{Xpp}$. These comparisons concern the
characteristic-zero points of this ring, including points that
require rescaling to an integral representation. On these points,
isomorphism over $\ker\eta$ is equivalent to the prescribed
diagonal torus orbit, and the native tangent comparison sends
rigidity to the formula $D(z)=c\,j\,\xi(z)$ on a homogeneous
element of weight $j$ for every point derivation $D$.

Apply $\mathrm{ThmS6\_s}$ to this universal pair. It gives the
finiteness just used and $B'_s(K,\iota)$ for every such type-I
point. Here $B'_s$ means all three exact requirements of
\texttt{BPrime\_s}, including irreducibility on open subgroups
of the Gaussian kernel, irreducibility on the intersection
with the full cyclotomic kernel, and $\mathrm{AutOnAlt}$ with
sign $-1$ for the pulled-back matrix representation. Hence
$\mathrm{HostRigid\_s}$ applies to each point with these same
parameters. Every two-sided point is separated by the
given $\phi$. The degree-zero element testing this is exactly
\[
 a=\operatorname{tr}u(\phi)-\operatorname{tr}r_z(\phi).
\]
Each diagonal entry has weight zero, the origin evaluates this
difference to zero, and $\xi(a)$ is its stated trace difference.

For precision about the algebraic step, if $d_{-1}=0$, the
injection of the next paragraph has zero image for every $b$,
so take $V=0$ and the empty list. Otherwise pad the homogeneous
relations by zero to exactly $d_{-1}+1,d_{-2},d_1-1,d_2-1,d_0$ equations of the respective weights $-1,-2,1,2,0$, as permitted by the five inequalities.
Nonpositive tangent vanishing means the respective linear forms
span the cotangent parts after inverting a uniformizer, by finite
dimensional linear duality. The weight-zero quotient is finite by
the marked theorem. The diagonal point comparison forces a fixed
point killing both signs to be the origin. The rigidity formula makes
two derivations at a two-sided point linearly dependent, since each
is a scalar multiple of the Euler derivation. The trace verifies the
last nonvanishing hypothesis. In making this identification one uses
the free weighted graded-algebra presentation with its kernel cut
out by the specified homogeneous equations and the original origin.
SP6, with valuation $e_O v_\ell$, bounds its finite weight-minus-one
cotangent length $V$ by $2 e_O\sum_i v_\ell(\xi_i(a))$. One may select
one representative per orbit from the finite cover and apply SP6
to the thinned cover; the estimate holds for every cover. The above
orbit comparison yields the required matrix nonisomorphism.

Finally use the native integral coordinate comparison: for each
$0\ne b\in O$ there is an additive injection
$F.\mathrm{ShaO}(b)\to(O/(b))^{W.\sigma}$ supported on weight
$-1$, annihilating every linearized defining relation. Choosing
arbitrary integral lifts of its coordinate values gives derivations
of the free weighted algebra at the origin; their reductions kill
the weight-minus-one equations. SP1 therefore bounds the length
of this module by the same $V$. The factor $e_O$ enters the upper
valuation bound exactly once. This derives the two displayed estimates
from the stated native model comparisons.
% END CVI5

% BEGIN CVI6
\paragraph{CVI6. Uniform exponent and the literal CountV bound.}
\label{claim:CVI6}
Assume the host inputs in the strengthened source signature and the
native marked-model comparisons of CVI5 for the frames of CVI2.
The resulting estimate is the full original $\mathrm{CountV}$.
Indeed choose the Gaussian field first. $\mathrm{HostNumber\_s}$
then gives one $C\in\mathbb N$ before any $\Omega,\ell,c,N,e,\psi$,
frame, coefficient embedding or marked element. For each of the
later instances choose $\iota$ and take the frame of CVI2, its
$g_1$ from $\mathrm{ThmS6\_s}$, and the $\phi$ of CVI5. This $\phi$
is fixed before every finite collection of points. The same
$\mathrm{HostNumber\_s}$ upper bound therefore applies to the
particular irredundant list of CVI5. With $X=\ell N e$ it says
\[
 \Bigl(\sum_i v_\ell(\operatorname{tr}r_i(\phi)-
                  \operatorname{tr}r_z(\phi))\Bigr)\log\ell\le X^C .
\]
VLT3 gives $\mathrm{Vlen}\le V<\infty$ at these original
frame lengths. Its normalization is division by the positive
integer $e_O$, so
\[
 \mathrm{Vn}\log\ell\le(V/e_O)\log\ell
 \le 2\Bigl(\sum_i v_\ell(\operatorname{tr}r_i(\phi)-
                         \operatorname{tr}r_z(\phi))\Bigr)\log\ell
 \le 2X^C\le X^{C+1}.
\]
The last inequality uses $X\ge2$, which holds for every permitted
prime, positive level and positive exponent. The same $C+1$ works
for every fixed target character $\chi$, including every $h$ with
$h.\chi=\chi$ and every $t$ in its original scope. CVI1 uses the
weak side $R_1\le2^{c_2(\psi)}$; the separate
$\mathrm{CountVs}$ concerns the strict reverse inequality.
% END CVI6

\paragraph{Boundary of this reconstruction.}
The strengthened formal implication takes exactly
$\mathrm{ThmS6\_s}$, the two algebraic lemmas, local finiteness,
local Tate duality, local Euler characteristic, global finiteness,
global Euler characteristic, normalized Poitou--Tate,
$\mathrm{HostVanE}$, $\mathrm{HostXpp}$,
$\mathrm{HostRigid\_s}$ and $\mathrm{HostNumber\_s}$.
The selected graph separately tracks the frame-length work;
the formal implication constructs its requisite native deformation
objects internally. CVI1--4 establish arithmetic, class-group and
cohomological components; VLT1--3 give the actual length comparison.
CVI5--6 explicitly retain the native marked universal pair,
its adapted presentation and equations, the point, tangent and
orbit comparisons, and the integral Selmer coordinate injection.
The pinned source supplies these in the coordinate proof of
\texttt{abstract\_count}. The readable specialization of the
previous general and one-character presentations to all these
native comparison statements is still outstanding here.

\endgroup
